\documentclass[12pt]{article}
\usepackage[utf8]{inputenc}
\usepackage[margin=0.75in]{geometry}
\usepackage{amsmath}
\usepackage{amssymb}
\usepackage{amsthm}
\usepackage{graphicx} % Required for inserting images
\usepackage{array}

\title{Norm Emergence Report: Resistance}
\author{Aaron Avram}
\date{September 13th 2026}

\begin{document}
\theoremstyle{plain}
\newtheorem{theorem}{Theorem}
\newtheorem{proposition}{Proposition}
\newtheorem{lemma}[theorem]{Lemma}
\newtheorem{corollary}[theorem]{Corollary}
\newtheorem{claim}[theorem]{Claim}

\theoremstyle{definition}
\newtheorem{definition}{Definition}
\newtheorem{example}{Example}
\newtheorem{exercise}[theorem]{Exercise}

\theoremstyle{remark}
\newtheorem{remark}[theorem]{Remark}
\newtheorem{note}[theorem]{Note}

\maketitle

\section{Introduction}
Overall, the student report has some useful ideas which can be added fairly quickly to the current model.
The report correctly identifies that the gossip mechanism is central to norm formation or subgroup resistance.
However, they treat it as the source of these behaviours and not as an emergent structure in a system predisposed to
them. Most of the decisions regarding what to keep or remove from the report follow from this observation.

\section{Theoretical Grounding: Resistance as an Alternative Social Reward Economy}

\subsection{What We Want to Model}
In the paper \textit{Social Reward Economies}, resistance is not modelled as a product
of generic deviance or systemic disruptions. Instead, it is viewed as the emergence of an alternate reward economy.
Specifically, when a subset $G_r \subseteq G$ of the set of total actors views the dominant schema $\sigma_d$ as detrimental relative to an alternative $\sigma_r$,
its members can form an alternative economy that rewards behaviours aligned with $\sigma_r$ instead of $\sigma_d$.
Concretely, we will view this as a gap in rewards on a subset $S_r$ of the total set of states $S_r$ where $\sigma_r$ generates significantly better rewards
for $G_r$ then $\sigma_d$. Resistance
is therefore not a property of individual agents but of a newly formed competing reward structure.
The goal of our simulation is to reproduce this formation process from end-to-end and its three possible outcomes.
These are described in the paper as (Eqs. (1)-(2) in section 6.2.2):
\begin{enumerate}
    \item[1.] $G_r$'s weighted feedback dominates $G \setminus G_r$'s for both schema which implies the alternative schema overturns
    the old norm (this is successful resistance / norm change, exemplified by \#MeToo in Section 6.2).
    \item[2.] The two economies' weighted feedback is comparable, this causes polarization: the economies coexist and
    compete without either dominating (Section 7.2.).
    \item[3.] $G \setminus G_r$'s feedback dominates implies $G_r$ fails to change the norm and remains a marginalized subgroup
    (Section 6.2.2, end).
\end{enumerate}
These are the target phenomena and are qualitatively different from the current simulation
which produces a single opinion leader and a single converged norm (the "cohesive social reward economy"
case of Section 7.1) with no mechanism by which a second, self-sustaining economy could ever appear. 

\subsection{Key Ideas We Want to Test}
To approach implementing the above phenomena, we need to condense Section 6 into implementable mechanisms.
Three of the key ideas that matter most are:
\begin{itemize}
    \item \textbf{Misalignment must be discovered, not assumed.} Resistance begins when agents in $G_r$
    notice that $u_i(s, \sigma_r) \gg u_i(s, \sigma_d)$ for $s \in S_r$ which is the interpretation of Section 6.1 in the above setup.
    This is a claim about what agents can observe, not an intrinsic property of the system hence the simulation needs a channel
    through which an agent can sample its own payoff under $\sigma_r$ before it can defect from $\sigma_d$.
    \item \textbf{Feedback and reputation are group-relative, not global.} The paper's notation in Section 6.2.1
    gives $G_r$ its own feedback function $f^r_i$ and its own influence weights $w_j^r$, explicitly distinct
    from the $f_j$ and the wider group $G$. Therefore an agent's standing in $G$ is different from its standing in $G_r$,
    and these two values could diverge once $G_r$ is formed as a subgroup in the system.
    \item \textbf{How do the size of $G_r$ and feedback drive resistance.} Section 6.2.3
    makes this explicit: $|G_r|$ trades off against how strong and low-variance $G_r$'s internal feedback
    is relative to the noise from the rest of $G$. This gives us a concrete empirical target:
    measure equations (1) and (2) in Section 6.2.2 across variations of the subgroup size and the reward-gap
    magnitude and observe whether qualitative regimes actually emerge in the simulation.
\end{itemize}

\subsection{What This Means for the Implementation}
The current engine converges every agent's reputation estimate to a single shared value as $t \to \infty$
(this is discussed in Section 3.1 below). This is a valid implementation of a cohesive social reward economy,
but it does not allow for the possibility of resistance as modelled above. This is because
the current simulation has no mechanism by which alternative feedback signal, alternative reputational weighting or an alternative
schema exist simultaneously with a dominant norm.

To augment the current code to test the ideas in Section 2.2 above requires four major changes
\begin{enumerate}
    \item[1.] \textbf{Persistent personal-utility tracking.} Currently, an agent only acts under its leader's policy,
    and it overwrites its personal-utility estimates whenever it optimizes reputation or status. Thus an agent in $G_r$
    following the dominant norm can never discover the gap $u_i(s, \sigma_r) \gg u_i(s, \sigma_d)$ which is precisely
    what the paper identifies as the trigger for resistance. This is the most important part of the simulation to change.
    \item[2.] \textbf{Sampling $\sigma_r$.} If an agent simply imitates its current leader,
    they will never encounter an alternative schema in the first placed so misalignment can be tracked
    but never observed. Thus, followers must have some non-zero probability of deviating from the leader's policy.
    \item[3.] \textbf{Group-relative feedback and reputation.} The quantities $f^r_i$ vs. $f_j$ and $w^r_j$ vs. $w$
    described in Section 6.2.1 are not the same object computed twice, they are different quantities that a subgroup
    computes for itself. Therefore, the implementation needs a mechanism by which a subgroup's internal reputational
    structure can diverge from the group's. Moreover, Section 4.1 describes gossip networks and opinion leaders
    as \textit{emergent structural properties} of the reward economy rather than externally imposed ones. Hence the divergence
    in the gossip network should arise from the dynamics of the simulation, not be hard coded as two separate networks
    or two separate parameter sets.
    \item[4.] \textbf{A reward structure encoding misalignment.} The paper's setup gives four
    distinct payoff distributions $(\sigma_r/\sigma_d \times G_r / G \setminus G_r)$, the current reward models
    have no way of representing this and need to be extended.
\end{enumerate}

These four requirements guide the decisions of the rest of this document. Every addition
kept from the report either directly implements one of them or supports the implementation and every rejection is
a case where the report's proposal would satisfy one of the above requirements through an artifically injected structure
rather than an emergent one like the paper actually describes.

\section{Additions From the Report}
\subsection{Gossip}
The student report identifies gossip as the central mechanism that connects agents and would drive resistance,
which is clear from the effect of gossip on the current simulations.
However, I believe they misclassified its role in the emergence of a resistant subgroup.

Namely, the student report suggested a modified gossip algorithm such that the unhappy subgroup
has a separate gossip network on top of the collective network for the entire group.
I believe this should be a emergent from the difference in reward structure between the two groups and not artificially created.

This is precisely the third requirement from Section 2.3: the paper treats gossip networks and opinion leaders
as emergent from the underlying feedback dynamics described in Section 4.1 of the paper,
not as infrastructure hardwired by the modeller. A manually split network would let two economies
coexist by construction, which tells us nothing about whether the resistance mechanism actually produces a split on its own.

However, the student report correctly identified that in the current simulation, there is no way for such a gossip network to form.
This is because each gossip update is of the form:
\[ s_i(k, t+1) = \frac{1}{|\mathcal A_p(t)|}\sum_{j \in \mathcal A_p(t)}s_j(k, t) + (v_i(k, t) - v_i(k, t - 1)) \]
where $s_i(k, t)$ is $i$'s reputation estimate for $k$ at time $t$, $\mathcal A_p(t)$ is the set of participating actors
at time $t$ and $v_i(k, t)$ is $i$'s estimate of its personal benefit from $k$ at time $t$.

Now as $t \to \infty$ theoretically we find that
\[ s_i(k, t) \approx s_j(k, t) \]
as the consensus term is identical for all participants and $v_i(k, t)$ plateaus as $t \to \infty$
which implies $v_i(k, t + 1) - v_i(k, t) \to 0$.

This can be verified empirically for sufficiently large $t$ close to norm convergence. This drives
the adoption of the current leader's norm and prevents any resistance from forming.
This exemplifies the claim made in Section 2.3 that the current update rule can only produce a single cohesive reward economy.

\subsection{Similarity}
Another proposed addition from the student report can resolve the gossip mechanism
by naturally creating the separated gossip networks for the total group and unhappy subgroup.
This addition is the similarity score, defined for a state $s \in S$ as
\[ \operatorname{sim}_t(i, j, s) = \begin{cases}
    \operatorname{sim}_{t - 1}(i, j, s) \cdot (1 - \beta) + \beta \frac{1}{|K_s(t)|}\sum_{k \in K_s(t)}\Big(u_i(t, s, \pi_k(s)) - u_j(t, s, \pi_k(s))\Big)^2 & t > 0 \\
    0 & t = 0
\end{cases}
\]
where $K_s(t)$ is the set of actors acting in state $s$ at time $t$. if $K_s(t) = \varnothing$, 
then there is no update.
for a fixed rate constant $\beta \in (0, 1)$
then
\[ \operatorname{sim}_t(i, j) = \frac{1}{|S|}\sum_{s \in S}\operatorname{sim}_t(i, j, s) \]
then there is a fixed constant $C$ such that agents $i$ and $j$ could gossip at time $t$
only if
\[ \operatorname{sim}_t(i, j) < C \]
and an intuitive hypothesis would be that overtime unhappy agents will have a low score with each other
and a high score with other agents and this would naturally fragment the gossip network
into two subnetworks.

This is exactly what is needed in requirement (3) in Section 2.3. The emergence of a $G_r$-specific gossip network
and $G_r$-specific weights $w^r_j$ is now a testable outcome of the dynamics rather than
a hard-coded property of the network structure. I think this is the one piece of the report closest to the paper's
own account of how alternative reputation structures form.

Additionally, a utility matrix is computed at each time step of the simulation, so implementing
such a similarity score would require minimal tweaking to the current model.

\subsection{Modelling the Unhappy Agents}
The report proposes to model a subset of unhappy agents as a subgroup $G_r$ of the total agents $G$
such that on a subset of the states $S_r \subseteq S$ the dominant schema $\sigma_d$
is significantly inferior to the potential resistant schema $\sigma_r$, i.e.
\[ u_i(s, \sigma_r) >> u_i(s, \sigma_d) \quad \forall i \in G_r, \forall s \in S_r \]
and this is a reasonable gap to drive the formation of a resistant subgroup that fragments from the dominant group.

This is a direct implementation of the misalignment condiiton in Section 6.1 of the paper, as this
quantity has to exist before resistance can even be discovered.

However, the report did not address how agents in $G_r$ will discover the inbalance in personal utility
under the dominant schema. This is because in the current simulation, the personal utility estimate of an agent
is unchanged when the agent is optimizing status or reputation. Additionally, for an agent that has adopted the dominant
schema as a follower of the current leader, it only acts under the policy of the leader without the chance to sample from its own
personal utility weights which is how its personal utility estimate could rise above its reputation estimate to detach it from the group.

This is requirement (1) from Section 2.3, without it the above gap can exist mathematically in the reward table
but never be visible to the agent and hence can never produce any observable behaviour in the system.

To fix these two issues, the personal utility estimate of each agent must be kept current and their personal
utility weights must not be overriden by the leader weights.

For the agents to actually occassionally witness the schema $\sigma_r$ I propose an extension of the report
model by adding another rate $\mu_\text{{dev}}$ where $\theta(\mu_{\text{dev}})$ gives the probability
that a given agent, regardless of whether they optimize personal utility, reputation or status,
will instead choose to optimize personal utility at that given time step. This deviation mechanism
will allow unhappy agents to explore other schemas which in turn can faciliate the formation of a resistant subgroup.

This is requirement (2) from Section 2.3, the discovery channel needed for observing of the alternative schema.

Note that for $\mu_{\text{dev}}$ to have any affect the current model has to be tweaked, as at the moment
an agent $i$ that optimizes reputation copies its leader's weights and does not keep track of its previous personal utility weights.

Another implementation detail, if two norms emerge in the simulation multiple functions in the model will fail silently
as they assume a single norm. These include paper\_welfare which takes a single leader\_id and the leader metric
which takes a global argmax of follower counts. Therefore, migrating the code from a single norm simulation to
one with two will require a lot of attention to which functionality implicitly assumes a single norm. 

\subsection{Attention Allocation for Participant Rate}
The student paper proposes a new way to model the interaction $\mu_i$. Namely
\[ \{ \mu_i(s)\}_{s\in S} = \operatorname{argmax}_{\mu_i} \sum_{s \in S} \theta_\alpha(\mu_{\text{min}} + \mu_i(s))\cdot u_i(s, \pi_i(s)) \]
subject to the constraints
\[ \sum_{s \in S}(\mu_{\text{min}} + \mu_i(s)) \leq M  \quad 0 \leq \mu_i(s) \]
where the agent i's interaction rate becomes $\mu_{\text{total}} = \mu_{\text{min}} + \mu_i$
here $M$ is a budget on the total attention an actor $i$ can allocate to different situations and $\mu_{\text{min}} > 0$ to prevent degenerate behaviour.
Under this model it is possible for an agent to selectively focus their attention on states that benefit them.

Additionally $\theta_\alpha$ is a generalization of $\theta$ which is defined as
\[ \theta_\alpha(\mu) = 1 - \exp(-\alpha\mu) \]
although up to scaling of $\mu$ and $M$, $\alpha$ is redundant so we take it to be equal to 1.

This optimization problem has a closed form solution from the KKT conditions.

Specifically there emerges a partition of the states $F \sqcup P = S$ called
the free and pinned states
where for all $s \in P$
\[ \mu_i(s)^* = 0 \]
then for $s \in F$
\[ \mu_i(s)^* = \frac{M}{|F|} - \Big(\mu_{\text{min}} + \frac{|P|}{|F|}\mu_{min} \Big) + \Big(\log u_i(s, \pi_i(s)) - \frac{1}{|F|}\sum_{t \in F}\log u_i(t, \pi_i(t))\Big) \]
which will be positive.

then for $s \in G_r$ still following the dominant schema $\sigma_d$ if for a moment they deviate
through $\mu_{\text{dev}}$ and act on personal utility they will update $\mu_i(s)$
based on $\sigma_r$ and this will drive $\mu_i(s)$ up for $s \in G_r$ and down for $s \in G \setminus G_r$.
Hence this attention mechanism shoud amplify deviation in the subgroup into cohension of the subgroup.

I kept this part of the report, because it implements the second half of the paper's success condition in Section 6.2.3:
resistance is not just about the existence of a utility gap, but about $G_r$'s feedback becoming disproportionately strong and coherent.
An attention mechanism that lets deviating agents concentrate their interaction budget on $S_r$ is a natural way to generate this disproportionate feedback
without injecting it artifically in the system.

The report did not specify where this model for $\mu_i$ would be used in the simulation, but I propose
that it replaces the current model for $\mu_{i, a}$.

An intuitive hypothesis would be that as the system progresses and the resistant subgroup
forms, a given member $s \in G_r$ its free states will correspond roughly to $S_r$
and its pinned states to $S \setminus S_r$. And the opposite effect would be present for $G \setminus G_r$.
This would reinforce a fragmented gossip network when combined with the similarity neighborhood.

In terms of implementation, the report does not give any method for how this interaction rate would
be updated during simulation. However, I designed a method that mirrors the approach for the current actor interaction
rate optimization.

For each agent $i$, given some initialization of $\mu_{i, a}(s, 0)$ for all $s \in S$ that satisfies
\[ \sum_{s \in S}(\mu_{i, a}(s) + \mu_{\text{min}}) \leq M \quad \forall s \in S, \mu_{i, a}(s) \geq 0  \]
at time $t$ the above optimization problem is solved and $\mu_{i, a}(s, t)^*$ are found for $s \in S$.
Then there is a fixed rate constant $\eta(t) \in (0, 1)$ such that the rate is updated to
\[ \mu_{i, a}(s, t+1) = (1 - \eta(t))\mu_{i, a}(s, t) + \eta(t) \cdot \mu_{i,a}(s, t)^* \]
and this clearly satisfies the constraints as $\{ \mu_{i, a}(s, t) \}_{s \in S}$ and $\{ \mu_{i, a}(s, t)^* \}_{s \in S}$
both do.

$\eta(t)$ should satisfy the same constraints as the rate variables in the learning norms paper. Namely
\begin{align*}
    &\sum_{n = 0}^\infty \eta(t)^2 < \infty \\
    &\sum_{n = 0}^\infty \eta(t) = \infty
\end{align*}

Now to actually calculate $\mu_i(s)^*$ in practice $F, P$ have to be computed.
To do this we sort $S$ as $\{ s_1, \ldots, s_N \}$ in increasing order with respect to $u_i(s_j, \pi_i(s_j))$.
It is clear that $P$ is of the form $\{ s_j : 0 < j \leq k \}$ for some $k \in \{ 0, \ldots, N + 1 \}$
Hence we can define
\[ L_k = \sum_{j \leq k}u_i(s_j, \pi_i(s_j))\theta(\mu_{\text{min}}) + \sum_{j > k}u_i(s_j, \pi_i(s_j))\theta(\mu_{i, a}(s_j, t) + \mu_{\text{min}}) \]
then $P = \{ s_j : 0 < j \leq k \}$ for
\[ k = \operatorname{argmax} \{ L_k : 0 \leq k \leq N + 1 \} \]
for every agent this has worst case complexity $\mathcal O(|S|\log |S|)$
and since this is done for each agent, every participant interaction rate update would have worst case complexity
$\mathcal O(|G||S|\log |S|)$ which is reasonable for moderately sized simulations.

The only implementation issue would be that this attention mechanism requires some significant
changes to the current code and could easily cause silent failures. Hence the process of migrating
from the current actor interaction rate to this attention-based one has to be done carefully.

Note that I took $\mu_{\text{min}}$ as uniform across all states to simplify the optimization, although in the report they allow $\mu_{\text{min}}$ to vary across states
which is a future extension to the above model.

Additionally this could be extended even further to $\mu_{i, p}$, the only difference is that instead
of weighting each state term by an agents own reward $u_i(s, \pi_i(s))$ it would be better to weight it by
\[ H_i(s) = \frac{1}{|G|}\sum_{k \in G}u_i(s, \pi_k(s)) \]
then solve the optimization problem
\[ \{ \mu_{i, p}(s)\}_{s\in S} = \operatorname{argmax}_{\mu_i} \sum_{s \in S} \theta_\alpha(\mu_{\text{min}} + \mu_{i, p}(s))\cdot H_i(s) \]
subject to the constraints
\[ \sum_{s \in S}(\mu_{\text{min}} + \mu_{i, p}(s)) \leq M  \quad 0 \leq \mu_i(s) \]
and the closed form solution can be found using the same method as the previous problem only replacing $u_i(s, \pi_i(s))$
with $H_i(s)$.

\subsection{Reward Table Structure}
The report proposes a reward table structure to model $\sigma_r$ and $\sigma_d$ that is intuitive
and produces $S_r$ and $G_r$ as discussed above. Namely for states $s \in S \setminus S_r$
agent rewards can be sampled under one mean and standard deviation. However for $s \in S_r$,
actions associated to the resistant schema $\sigma_r$ will a reward distribution given by
\[ u_i(s, \sigma_r) \sim \mathcal N(m_r^{G_r}, {\sigma_r^{G_r}}^2) \quad \forall i \in G_r\]
and
\[ u_i(s, \sigma_r) \sim \mathcal N(m_r^{G}, {\sigma_r^{G}}^2) \quad \forall i \in G \setminus G_r\]
then for the dominant schema $\sigma d$:
\[ u_i(s, \sigma_d) \sim \mathcal N(m_d^{G_r}, {\sigma_d^{G_r}}^2) \quad \forall i \in G_r \]
and
\[ u_i(s, \sigma_d) \sim \mathcal N(m_d^{G}, {\sigma_d^{G}}^2) \quad \forall i \in G \setminus G_r\]
where the marginalized group is created by enforcing
\[ m_r^{G_r} >> m_d^{G_r}  \quad m_d^{G} >> m_r^{G} \]
this sort of reward table can be built on top of the current reward table structure in the simulation
without many changes.

I kept this section essentially as proposed because it is a direct, minimal implementation of requirement (4)
from Section 2.3 and the natural discretization of the reward model described in Section 6.2.1 of the paper.

\section{What Will Not be Used}

\subsection{Multiple Gossip Networks}
The report proposed a manually fragmented gossip network, which I have decided to scrap
in favour of using the above similarity neighborhoods to produce a fragmented gossip network
naturally. My approach is aligned with the social reward economies paper which treats
gossip networks as emergent (Section 4.1: gossip networks "can both stabilize and destabilize power structures"
as a consequence of ongoing reputational dynamics, not as a fixed architecture). Additionally, Section 7.2.1
of the paper describes polarized subgroups as "structurally insulated social systems" that become insulated
by their own internal feedback amplification. Hence, insulation should be an outcome of the dynamics in the paper,
not an input to them.

\subsection{Separate Gammas}
The model proposes several $\gamma$ values that modulate the weight of an agent $i$'s reputation in $G$ versus $G_r$
for $i \in G_r$ versus $i \in G \setminus G_r$. This is an artifical addition
as it indexes $\gamma$ by group membership to create a gap in reputation for different groups.
However, this difference is precisely what should emerge from the fragmentation of the gossip network mentioned above.
The paper is explicit that this gap is earned not assigned. Section 6.2.1 states: agents in $G_r$ "who strongly support the schema $\sigma_r$
and strongly comend $\sigma_d$ gain reputation within the subgroup, thereby acquiring higher influence weights $w_j^r$."
Hence $w_j^r$ is itself an output of the within-subgroup feedback process, not a parameter.
Selecting separate $\gamma$'s per group would bake in the very asymmetry the model is supposed to explain.

\subsection{Experiment Setup}
All of the experiment set up in the report makes sense, any parameters which I have removed above such as the different $\gamma$
values will be removed. However, one point where there the experiment should have a different setup is at initialization,
namely the current experiment seeds all agents with the dominant norm except those in the unhappy group $G_r$ which are given $\sigma_r$.
However this artifically forms the resistant subgroup and only studies whether or not it survives over time. This is a separate question
from whether or not this subgroup will form on its own. Before the changes to personal utility tracking from above are implemented,
it would have been impossible for the resistant subgroup to form and seeding the resistant subgroup with the resistant schema
would have been necessary for any behaviour to emerge. However, with my proposed changes to personal utility tracking and follower deviation,
the subgroup formation can happen naturally. This distinction matters to accurately test the claims of the paper. Section $6.1$
defines resistance as something that emerges once misalignment is discovered. Seeding $G_r$ with $\sigma_r$
from $t = 0$ answers a survival question (does an existing alternative economy persist against the dominant one)
rather than the formation question the report and section $6.1$ of the paper are discussing.
With requirements (1) and (2) from section 2.3 in place, seeding is no longer necessary and the initialization
should instead start with every agent under $\sigma_d$, letting $G_r$ (if it forms at all) be an outcome
we measure rather than a condition we impose.

\subsection{Parameter Variation in Experiments}
The report proposes that each of its experiment parameters are varied one at a time for each experiment run,
this is not optimal as it assumes the effect of each parameter is additive.
Since I have already cut down the number of parameters in the simulation I propose instead to run a grid of different parameter values.
Even if this does not mean all parameters are in the grid, even a experiment with a grid of 2 or 3 parameters
will give richer results than an experiment that only varies one. This also follows from Section 2.2's third point:
Section 6.2.3 frames resistance as a joint function of subgroup size $|G_r|$ and feedback, not either individually.
Thus, single parameter sweeps cannot actually test the paper's success condition. Only a joint grid
over size, reward-gap (and now, the deviation rate and similarity threshold) can show whether the two factors trade off against
each other in the way the paper predicts.

In terms of specific parameters, the fixed parameters the report proposed which will still be present are
\begin{enumerate}
    \item[1] N: the number of agents
    \item[2] S: the set of situations
    \item[3] M: the total attention budget
    \item[4] $\gamma$: the reputation scaling factor
    \item[5] $\kappa$: the status scaling factor
    \item[6] Normal distribution parameters for reward tables 
\end{enumerate}
and the variable parameters are
\begin{enumerate}
    \item[1] The distribution of $\mu_{i, a}(s)$.
    \item[2] The distribution of $\mu_{i, p}(s)$.
    \item[3] The distribution of $\mu_{\text{min}}(s)$ (if taken to be non-uniform) versus total utility for each situation.
\end{enumerate}
Additionally some parameters not mentioned in the report that would be added are
\begin{enumerate}
    \item[1] C: the similarity score threshold to enable gossip.
    \item[2] $\mu_{\text{dev}}$: the rate modulating the chance of follower deviation. 
\end{enumerate}

\end{document}